% Generated draft: requires mathematical and source review.
\documentclass[11pt]{article}
\usepackage[margin=25mm]{geometry}
\usepackage{fontspec}
\setmainfont{Latin Modern Roman}
\usepackage{amsmath,amssymb,mathtools,unicode-math}
\setmathfont{Latin Modern Math}
\usepackage{enumitem,needspace,xurl,hyperref}
\setlength{\parindent}{0pt}
\setlength{\parskip}{5pt}
\setlength{\emergencystretch}{3em}
\newcommand{\N}{\mathbb{N}}
\newcommand{\Z}{\mathbb{Z}}
\newcommand{\Q}{\mathbb{Q}}
\newcommand{\R}{\mathbb{R}}
\newcommand{\C}{\mathbb{C}}
\newcommand{\F}{\mathbb{F}}
\newcommand{\A}{\mathbb{A}}
\newcommand{\PP}{\mathbb{P}}
\newcommand{\E}{\mathbb{E}}
\newcommand{\eps}{\varepsilon}
\newcommand{\dd}{\,\mathrm{d}}
\begin{document}
\section*{UnsolvedMath ID: 10600035}
\textbf{Source problem code:} catalogue entry
\subsection{Definitions and mathematical statement}
Let \(F_K\) denote the source's Kauffman polynomial of a link \(K\), and \(P_L\) its HOMFLY--PT polynomial. The Jaeger expansion referred to in item 35 has the schematic form
\[
 F_K=\sum_{s\in\mathcal S(K)}w_s\,P_{K_s},
\]
where \(K_s\) are oriented, partially smoothed links associated with \(K\), and \(w_s\) are state weights. The source record supplies neither the actual state set and local smoothings nor their weights, variables or normalization; the formula above records only the stated structure.

The request is to generalize this Jaeger viewpoint to categorify the Kauffman polynomial, using chain complexes for the states derived from the Khovanov--Rozansky categorification of HOMFLY--PT. In a precise construction one must assemble those state complexes with the necessary grading data and differentials into a link-invariant homological object whose appropriately normalized graded Euler characteristic is \(F_K\). The local weights, compatibility maps and invariance data are not determined by a bare direct sum of state homologies, and they are not specified in the extract.

No particular coefficient ring, grading convention, link category, or reduced/unreduced version is supplied. The request concerns this state-based construction for the source's Kauffman polynomial; a fixed \(\mathrm{SO}(2N)\) specialization is not automatically the full two-variable problem. The supplied partial-progress assessment reports relevant specialized categorifications and a Jaeger progress report, while saying the full state-categorification program appears incomplete; these are source-qualified assessments, not independently verified status claims.
\subsection{Short English statement}
Lift Jaeger's HOMFLY--PT state expansion of the Kauffman polynomial to a homological construction using Khovanov--Rozansky complexes for its oriented states. The source's exact expansion and the assembly maps are missing. Reported specialized categorifications are retained as partial progress rather than automatically identified with this full request.
\subsection{Sources}
\begin{itemize}
\item Local UnsolvedMath registry, permanent ID \(10600035\), exact source code \texttt{catalogue entry}, virtual-knot problem-list item 35. The expansion is attributed in the statement to Jaeger, and the input categorification to Khovanov--Rozansky. The opening reference slot is empty; the title is truncated.
\item \emph{Unsolved Problems in Virtual Knot Theory and Combinatorial Knot Theory}, revised survey, dedicated problem-list item 35: \url{https://arxiv.org/abs/1409.2823}. Background: 60 top-level items, compound items retained intact; source-TeX extraction, accessed 2026-07-29.
\item Supplied incomplete references: ``Virtual crossings, convolutions and a categorification of the SO(2N) Kauffman polynomial'' (2007), arXiv:0701.xxxx; ``On Jaeger's HOMFLY-PT expansions, branching rules and link homology: a progress report'' (2013), arXiv:1309.xxxx. The xxxx identifiers are placeholders and have not been completed.
\item Registry status partially\_solved and classification PARTIAL-PROGRESS report specialized categorifications and progress toward the Jaeger-state approach, without establishing completion of the full request. Source evidence NEEDS\_REVIEW calls for a current primary-literature audit.
\item Dataset attribution: Ulam AI / UnsolvedMath Contributors, CC BY 4.0. Source rights NEEDS\_REVIEW: publicly accessible author/publisher copy, redistribution terms requiring review.
\end{itemize}
\end{document}
